\section{The Compositional Refinement Pass}
\label{sec:algorithm}

Our pass takes as input a populated \texttt{SpecificationCache}, a call graph over the same methods, and a map from method signature to the abstract syntax tree (AST) of each \texttt{MethodDeclaration}. We mutate the cache in place.
Figure~\ref{fig:pipeline} sketches the data flow.

\begin{figure}[!t]
\centering
\begin{tikzpicture}[
  font=\scriptsize\sffamily,
  node distance=3mm and 5mm,
  artefact/.style={
    rectangle, rounded corners=1.5pt, draw=black!55, fill=black!4,
    minimum height=4.5mm, minimum width=18mm, inner sep=2pt, align=center
  },
  pass/.style={
    rectangle, draw=black!80, fill=blue!8,
    minimum height=4.5mm, minimum width=22mm, inner sep=2pt, align=center
  },
  predecessor/.style={
    rectangle, draw=black!60, dashed, fill=gray!6,
    minimum height=4.5mm, minimum width=22mm, inner sep=2pt, align=center
  },
  >={Stealth[length=1.5mm,width=1.3mm]},
]
  % Row 1: single-pass predecessor
  \node[artefact] (src) {Source jar};
  \node[predecessor, right=of src] (inferrer) {Single-pass\\inferrer (predecessor)};
  \node[artefact, right=of inferrer] (cache) {\texttt{SpecCache}};
  % Row 2: this paper
  \node[pass, below=4mm of inferrer] (driver) {\texttt{CompositionalAnalyzer}\\\scriptsize SCC driver};
  \node[pass, below=of driver] (perMethod) {Per-method WP walk};
  \node[artefact, right=of perMethod] (refined) {Refined\\\texttt{SpecCache}};
  \draw[->] (src) -- (inferrer);
  \draw[->] (inferrer) -- (cache);
  \draw[->] (cache.south) |- (driver.east);
  \draw[->] (driver) -- (perMethod);
  \draw[->] (perMethod) -- (refined);
  \draw[->, dotted] (perMethod.west) .. controls +(left:6mm) and +(left:6mm) .. (driver.west) node[midway, left=1pt, font=\tiny, align=center] {callee\\specs};
\end{tikzpicture}
\caption{The compositional refinement pass. The single-pass predecessor (dashed) populates a per-method \texttt{SpecCache}; \texttt{CompositionalAnalyzer} iterates strongly-connected components of the call graph in reverse topological order, and for each method walks the body backwards (WP) refining the cached spec against the now-refined callee preconditions. The cache is mutated in place.}
\label{fig:pipeline}
\end{figure}

\subsection{Driver}

Our driver iterates strongly-connected components (SCCs) of the call graph in reverse topological order, so callees are refined before callers. For each SCC:

\begin{itemize}\setlength\itemsep{0.2em}
  \item \emph{Singleton SCC (no recursion):} refine the single method
        once.
  \item \emph{Cyclic SCC (mutual recursion):} refine each method in
        the SCC; repeat until no method gains a new clause or until
        a configurable fixpoint cap (default: 16 iterations) is
        reached. The cap exists to bound runtime on pathological
        graphs; we log a warning when it fires.
\end{itemize}

\subsection{Per-method refinement}

For each method \texttt{m}, the refiner walks every
\texttt{MethodCallExpr} in \texttt{m}'s body and, for each call site,
performs the following sequence (Algorithm~\ref{alg:refine}):

\begin{algorithm}[t]
\caption{Compositional refinement.}
\label{alg:refine}
\small
\begin{algorithmic}[1]
\Require populated cache $C$, call graph $G$, declaration map $D$
\Ensure $C$ refined in place
\State $\textit{scc} \gets$ \Call{ReverseTopologicalSCCs}{$G$}
\ForAll{$S \in \textit{scc}$}
  \Repeat
    \State $\textit{changed} \gets \mathbf{false}$
    \ForAll{$m \in S$}
      \State $\textit{spec} \gets C(m)$
      \ForAll{$\textit{call} \in $ \Call{CallSites}{$D(m)$}}
        \State $g \gets$ \Call{EnclosingGuards}{$D(m), \textit{call}$}
        \If{\Call{SideEffecting}{$g$} \textbf{or} \Call{SideEffecting}{$\textit{call}.\text{args}$}}
          \State \textbf{continue} \Comment{reject unsound substitutions}
        \EndIf
        \State $\textit{candidates} \gets$ \Call{ResolvePolymorphic}{$\textit{call}, G$}
        \State $\textit{lifted} \gets \emptyset$
        \ForAll{$c \in \textit{candidates}$}
          \State $\sigma \gets$ \Call{SubstituteArgs}{$C(c).\text{pre}, \textit{call}.\text{args}$}
          \State $\textit{lifted} \gets \textit{lifted} \cup \{ g \Rightarrow \sigma \}$
        \EndFor
        \If{$|\textit{candidates}| > 1$}
          \State $\textit{spec}.\text{pre} \gets \textit{spec}.\text{pre} \cup \{\bigvee \textit{lifted}\}$
        \Else
          \State $\textit{spec}.\text{pre} \gets \textit{spec}.\text{pre} \cup \textit{lifted}$
        \EndIf
        \If{\Call{Terminates}{$C(c)$}}
          \State $\textit{spec}.\text{post} \gets \textit{spec}.\text{post} \cup \text{(lifted postconditions)}$
        \EndIf
      \EndFor
      \If{$\textit{spec} \ne C(m)$} $C(m) \gets \textit{spec}$; $\textit{changed} \gets \mathbf{true}$
      \EndIf
    \EndFor
  \Until{$\neg \textit{changed}$ \textbf{or} cap reached}
\EndFor
\end{algorithmic}
\end{algorithm}

\begin{enumerate}\setlength\itemsep{0.2em}
  \item \textbf{Resolve the callee.} Match the call's simple name
        against keys in the cache to find a candidate \texttt{n}. A
        full type-resolving harness is not required for the present
        measurement; mismatches caused by overloading are documented
        as a measurement limitation (Section~\ref{sec:threats}).
  \item \textbf{Substitute arguments.} For each precondition of
        \texttt{n}, replace each formal parameter name with the actual
        argument expression. Substitution is regex-based with word
        boundaries so a formal \texttt{s} does not clobber a substring
        \texttt{this.size}. Side-effecting actuals
        (\texttt{++}, \texttt{--}) reject the substitution as unsound.
  \item \textbf{Lift through enclosing guards.} Walk the AST upward
        from the call expression and accumulate the conjunction of
        any \texttt{if}-condition or ternary-condition whose body
        contains the call. If guards are present, the substituted
        precondition becomes \texttt{(guard) ==> (substituted)}.
        Calls in the \texttt{else} branch contribute \texttt{!(guard)}.
        Side-effecting guards are rejected to preserve soundness.
  \item \textbf{Detect polymorphic dispatch.} If the resolved callee
        has overriding implementations in subclasses (looked up via
        the call graph's class-hierarchy index), collect every
        candidate's substituted preconditions and emit the disjunction
        \texttt{(p\_1 || p\_2 || \dots\ || p\_k)}. The disjunction is
        sound: at least one candidate's contract must hold for the
        dispatch to be safe.
  \item \textbf{Conjoin into the caller's set.} Add the lifted clause
        to \texttt{m}'s precondition list if it is not already
        present.
\end{enumerate}

\subsection{Termination-aware postcondition propagation}

The pass also propagates the callee's \emph{postconditions} into the
caller's set, gated on a termination flag the inferrer maintains for
each method. A postcondition of a non-terminating callee describes a
state that is established only \emph{if} the callee returns, so
propagating it unconditionally would be unsound when the callee may
not terminate. The pass therefore propagates \texttt{n}'s
postconditions into \texttt{m} only when \texttt{n.terminates()} is
true. The polymorphic-dispatch case is handled identically to
preconditions, with the disjunction over candidates.

\subsection{What the algorithm is and is not}

The algorithm is a \emph{lifting} step over an existing single-pass
result, not a complete WP transformer. It does not:

\begin{itemize}\setlength\itemsep{0.2em}
  \item compute the weakest precondition through arithmetic or array
        statements (it accepts the inferrer's clauses verbatim);
  \item handle aliasing (the inferrer's substitution is syntactic);
  \item lift through loops (the inferrer's loop-invariant heuristic
        handles this in the isolated pass);
  \item discharge clauses against a satisfiability-modulo-theories (SMT) solver (the inferrer's
        downstream OpenJML pipeline handles this).
\end{itemize}

It does enough --- branch-guard lifting and polymorphic-dispatch
disjunction --- to expose the structural classes of additions that a
single-pass propagation strategy cannot produce. Section~\ref{sec:empirical}
measures whether those additions are empirically substantial on real
code.
